\begin{frame}{FAQ: 은닉 상태 크기와 인코더-디코더}
  \begin{block}{Q. 은닉 상태의 크기는 어떻게 정하나요?}
    ``지금까지의 맥락을 얼마나 \textbf{풍부하게 요약}할 수 있는가''를
    결정하는 하이퍼파라미터 --- Week 9의 은닉층 뉴런 수와 같은
    종류의 선택.
    \begin{itemize}
      \item 너무 작으면 \textbf{긴 문맥을 다 담지 못}하고
      \item 너무 크면 파라미터 증가 $\to$ \textbf{과적합 · 계산비용}
        문제
    \end{itemize}
  \end{block}
  \vskip0.6em
  \begin{block}{Q. 입력과 출력이 다른 경우(번역)는?}
    ``입력을 끝까지 읽고 \textbf{마지막 은닉 상태 하나}로부터 새
    시퀀스를 순차 생성''하는 \kb{인코더-디코더(encoder-decoder)}
    구조를 쓰면 가능 --- 단, 긴 입력을 하나의 벡터로 압축하면
    \textbf{정보 손실}이 생김.
    \kq{Week 12의 Attention이 정확히 이 한계를 보완하기 위해
    출발했다} (Bahdanau et al., 2014).
  \end{block}
\end{frame}
